\subsection{Exact accelerated proximal iterates}

Apply accelerated proximal point directly to the orthant objective
$F_{\bar\rho}^+$. With $\kappa_{\mathrm A}=1-2\alpha$, define
\begin{equation}
 \vartheta_{\mathrm A}:=
 \frac{\alpha}{\alpha+\kappa_{\mathrm A}},
 \qquad
 \beta_{\mathrm A}:=
 \frac{1-\sqrt{\vartheta_{\mathrm A}}}
      {1+\sqrt{\vartheta_{\mathrm A}}},
 \label{eq:app-momentum}
\end{equation}
and initialize $x_{-1}=x_0=0$.  For $t\ge1$ let
\begin{align}
 y_{t-1}&=(1+\beta_{\mathrm A})x_{t-1}-\beta_{\mathrm A} x_{t-2},
 \label{eq:app-center}\\
 x_t&=\arg\min_{x\ge0}
 \left\{F_{\bar\rho}(x)
 +\frac{\kappa_{\mathrm A}}{2}\norm{x-y_{t-1}}_2^2\right\}.
 \label{eq:app-step}
\end{align}
Set
\[
 A_t:=\supp(x_t),
 \qquad V_t:=\vol(A_t),
 \qquad p_t:=p(x_t).
\]
The relevant nonmonotonicity is not the number of new coordinates.  It is the
mass of coordinates present two iterates ago and absent now:
\begin{equation}
 d_t
 :=p_{A_{t-2}\setminus A_t}(x_{t-2})
 =\sum_{i\in A_{t-2}\setminus A_t}w_i x_{t-2,i}.
 \label{eq:dropout-mass}
\end{equation}

\subsection{One-step active-volume inequality}

\begin{lemma}[Active KKT sum; \ProvedStatus]
\label{lem:active-kkt-sum}
Every exact step~\eqref{eq:app-step} satisfies
\begin{equation}
 (\alpha+\kappa_{\mathrm A})p_t+\alpha\bar\rho V_t
 \le \alpha+\kappa_{\mathrm A} p_{A_t}(y_{t-1}).
 \label{eq:active-kkt-sum}
\end{equation}
\end{lemma}

\begin{proof}
On $A_t$, the proximal KKT equations are
\[
 \bigl(Qx_t-b+\alpha\bar\rho w
       +\kappa_{\mathrm A}(x_t-y_{t-1})\bigr)_i=0.
\]
Because $x_t$ is zero off $A_t$ and $Q$ has nonpositive off-diagonal
entries,
\[
 \sum_{i\in A_t}w_i(Qx_t)_i
 \ge w^TQx_t=\alpha p_t.
\]
Also $\sum_{i\in A_t}w_i b_i\le w^Tb=\alpha$.  Weighted summation of the
active equations gives~\eqref{eq:active-kkt-sum}.
\end{proof}

The restricted center mass obeys
\begin{align}
 p_{A_t}(y_{t-1})
 &=(1+\beta_{\mathrm A})p_{A_t}(x_{t-1})
   -\beta_{\mathrm A} p_{A_t}(x_{t-2}) \notag\\
 &\le (1+\beta_{\mathrm A})p_{t-1}-\beta_{\mathrm A} p_{t-2}+\beta_{\mathrm A} d_t.
 \label{eq:center-mass-bound}
\end{align}
Support growth is free in this inequality; only disappearance of old mass
creates the final positive term.

\subsection{Cumulative identity}

\begin{theorem}[Cumulative active volume; \ProvedStatus]
\label{thm:cumulative-volume}
For the exact iterates~\eqref{eq:app-center}--\eqref{eq:app-step},
\begin{equation}
 \alpha\bar\rho\sum_{t=1}^T V_t
 +\alpha\sum_{t=1}^T p_t
 \le
 \alpha T+\kappa_{\mathrm A}(\beta_{\mathrm A} p_{T-1}-p_T)
 +\kappa_{\mathrm A}\beta_{\mathrm A}\sum_{t=1}^T d_t.
 \label{eq:cumulative-volume-identity}
\end{equation}
Define the nonnegative deactivation flux
\begin{equation}
 \Phi_T
 :=\pos{\beta_{\mathrm A} p_{T-1}-p_T}
   +\beta_{\mathrm A}\sum_{t=1}^T d_t.
 \label{eq:flux}
\end{equation}
Then
\begin{equation}
 \sum_{t=1}^T V_t
 \le \frac{T}{\bar\rho}
     +\frac{\kappa_{\mathrm A}\Phi_T}{\alpha\bar\rho}.
 \label{eq:volume-via-flux}
\end{equation}
\end{theorem}

\begin{proof}
Substitute~\eqref{eq:center-mass-bound} into
~\eqref{eq:active-kkt-sum}, separate $\alpha p_t$ from
$(\alpha+\kappa_{\mathrm A})p_t$, and sum over $t$.  The momentum terms telescope:
\[
 \sum_{t=1}^T
 \bigl((1+\beta_{\mathrm A})p_{t-1}-\beta_{\mathrm A} p_{t-2}-p_t\bigr)
 =\beta_{\mathrm A} p_{T-1}-p_T,
\]
because $p_{-1}=p_0=0$.  This proves
~\eqref{eq:cumulative-volume-identity}.  Replace the endpoint by its positive
part and discard the nonnegative term $\alpha\sum_t p_t$ to obtain
~\eqref{eq:volume-via-flux}.
\end{proof}

For an approximate inner solve, let
\begin{equation}
 E_t:=\sum_{i\in A_t}w_i
 \left|
   \bigl(Qx_t-b+\alpha\bar\rho w
          +\kappa_{\mathrm A}(x_t-y_{t-1})\bigr)_i
 \right|.
 \label{eq:inner-kkt-error}
\end{equation}
The right side of~\eqref{eq:volume-via-flux} gains
$\sum_tE_t/(\alpha\bar\rho)$.

Since $Q+\kappa_{\mathrm A} I$ has constant condition number, a positive-residual local
coordinate method solves each proximal problem in $\softO(V_t)$ work.  Thus
\begin{equation}
 \Work_{\mathrm{warm}}
 =\softO\!\left(
   \frac{T}{\rho}
   +\frac{\Phi_T}{\alpha\rho}
   +\frac{\sum_tE_t}{\alpha\rho}
 \right).
 \label{eq:warm-work-flux}
\end{equation}
Combining with~\eqref{eq:tail-target} and
$T=\softO(1/\sqrt\alpha)$ gives the observable conditional bound
\begin{equation}
 \Work_{\mathrm{hybrid}}
 =\softO\!\left(
   \frac1{\rho\sqrt\alpha}
   +\frac{\Phi_T}{\alpha\rho}
   +\frac{\sum_tE_t}{\alpha\rho}
 \right).
 \label{eq:hybrid-flux-bound}
\end{equation}

\subsection{The flux conjecture and transformed wave}

The tempting final lemma was
\begin{conjecture}[Deactivation-flux conjecture; now rejected]
\label{conj:flux}
For the exact accelerated Euclidean proximal iterates above, uniformly over
graphs,
\begin{equation}
 \Phi_T=O(\alpha T).
 \label{eq:flux-conjecture}
\end{equation}
\end{conjecture}
If true, it would turn~\eqref{eq:hybrid-flux-bound} into the target
~\eqref{eq:target-work}.  It is strictly sharper than counting face changes:
monotone support growth has $d_t=0$ regardless of how many vertices are
added, and a dropout is charged by its value rather than its degree.

The correct dynamic variable reveals why the conjecture is fragile.  Let
\[
 r_0:=1-\sqrt{\frac{\alpha}{\alpha+\kappa_{\mathrm A}}},
 \qquad
 z_t:=x_t-r_0x_{t-1}.
\]
Writing $\lambda_t\ge0$ for the orthant KKT multiplier, the inertial system
factors as
\begin{equation}
 (\alpha+\kappa_{\mathrm A})(z_t-r_0z_{t-1})
 +(Q-\alpha I)x_t
 =b-\alpha\bar\rho w+\lambda_t.
 \label{eq:transformed-wave}
\end{equation}
Both the endpoint term and a two-step dropout are controlled by negative
parts of the transformed momentum.  Consequently, the conjecture is, up to
constants, a bound on the cumulative weighted negative wave.  Branching
graphs can amplify that wave in degree-weighted $\ell_1$ even though the
accelerated method remains stable in objective and Euclidean norm.  The exact
mechanism is developed in \cref{sec:tree-wave}.
